\documentclass[10pt]{article}

% ---------- Document Identity (update these only) ----------
\newcommand{\DocStage}{}
\newcommand{\DocTitle}{Your Road to Tier One SOF}
\newcommand{\ProgramName}{182-Week Civilian-to-Selection Readiness Pipeline}
\newcommand{\ProgramSubtitle}{}

% ---------- Page ----------
\usepackage[legalpaper,left=0.5in,right=0.5in,top=0.6in,bottom=0.45in]{geometry}

% ---------- Typography ----------
\usepackage[T1]{fontenc}
\usepackage[utf8]{inputenc}
\usepackage{libertinus}
\usepackage{microtype}
\usepackage{parskip}
\usepackage{ragged2e}
\usepackage{textcomp}
\AtBeginDocument{\color{black}}
\AtBeginDocument{\pagecolor{white}}
\AtBeginDocument{\justifying}

% ---------- Landscape pages ----------
\usepackage{pdflscape}

% ---------- Color ----------
\usepackage[table]{xcolor}
\definecolor{StageBlue}{HTML}{1F4E79}
\definecolor{StageBlueDark}{HTML}{163A5A}
\definecolor{LightGray}{HTML}{F2F4F7}
\definecolor{MidGray}{HTML}{D9DEE5}
\definecolor{RuleGray}{HTML}{C7CED8}
\definecolor{HeaderInk}{HTML}{000000}
\definecolor{Ink}{HTML}{000000}

% ---------- Utility ----------
\usepackage{enumitem}
\sloppy
\setlength{\emergencystretch}{2em}

% ---------- Boxes / UI ----------
\usepackage[most]{tcolorbox}

\tcbset{
  charterTitle/.style={
    enhanced,
    colback=StageBlue!4,
    colframe=StageBlue,
    boxrule=1.25pt,
    arc=3pt,
    left=16pt,right=16pt,top=14pt,bottom=14pt,
    borderline north={3.2pt}{0pt}{StageBlue},
    borderline west={4pt}{0pt}{StageBlue},
    borderline south={1.25pt}{0pt}{StageBlue},
    drop shadow={black!25!white},
  },
  charterRule/.style={
    enhanced,
    colback=LightGray,
    colframe=RuleGray,
    boxrule=0.85pt,
    arc=3pt,
    left=12pt,right=12pt,top=10pt,bottom=10pt,
    borderline west={3pt}{0pt}{StageBlue},
  },
  charterBadge/.style={
    enhanced,
    colback=StageBlue,
    coltext=white,
    colframe=StageBlueDark,
    boxrule=0.6pt,
    arc=2pt,
    left=6pt,right=6pt,top=3pt,bottom=3pt,
  }
}
\newtcbox{\Badge}{nobeforeafter,charterBadge}

% ---------- Lists ----------
\setlist[itemize]{leftmargin=1.5em, itemsep=0.25em, topsep=0.25em}
\setlist[enumerate]{leftmargin=1.8em, itemsep=0.25em, topsep=0.25em}

% ---------- TOC ----------
\usepackage{tocloft}
\setcounter{tocdepth}{3}
\setlength{\cftbeforesecskip}{3pt}
\setlength{\cftbeforesubsecskip}{2pt}
\setlength{\cftbeforesubsubsecskip}{0pt}
\renewcommand{\cftsecfont}{\bfseries}
\renewcommand{\cftsecpagefont}{\bfseries}
\renewcommand{\cftsubsecfont}{\normalfont}
\renewcommand{\cftsubsecpagefont}{\normalfont}
\renewcommand{\cftsubsubsecfont}{\normalfont}
\renewcommand{\cftsubsubsecpagefont}{\normalfont}
\setlength{\cftsecindent}{0pt}
\setlength{\cftsubsecindent}{1.25em}
\setlength{\cftsubsubsecindent}{2.8em}
\setlength{\cftsecnumwidth}{2.1em}
\setlength{\cftsubsecnumwidth}{2.8em}
\setlength{\cftsubsubsecnumwidth}{3.6em}

% Remove the built-in TOC heading so only the custom heading is shown
\makeatletter
\renewcommand{\tableofcontents}{\@starttoc{toc}}
\makeatother

% ---------- Header/Footer: page number only ----------
\usepackage{fancyhdr}
\pagestyle{fancy}
\fancyhf{}
\renewcommand{\headrulewidth}{0pt}
\renewcommand{\footrulewidth}{0pt}
\fancyfoot[C]{\thepage}

% ---------- Section formatting ----------
\usepackage{titlesec}

\titleformat{\section}
  {\Large\bfseries\color{Ink}}
  {\thesection}{0.6em}{}
  [\vspace{0.25em}\textcolor{StageBlue}{\rule{\linewidth}{0.8pt}}]

\titleformat{\subsection}
  {\large\bfseries\color{Ink}}
  {\thesubsection}{0.6em}{}

\titleformat{\subsubsection}
  {\normalsize\bfseries\color{Ink}}
  {\thesubsubsection}{0.6em}{}

\titlespacing*{\section}{0pt}{0.95em}{0.35em}
\titlespacing*{\subsection}{0pt}{0.65em}{0.25em}
\titlespacing*{\subsubsection}{0pt}{0.55em}{0.2em}

% ---------- Table engine ----------
\usepackage{tabularray}
\UseTblrLibrary{booktabs}

% ---------- tabularray: GLOBAL table treatment ----------
\SetTblrInner{
  cells = {fg=black, bg=white, valign=t},
}

% ---------- Header helpers ----------
\newcommand{\Hdr}[1]{\shortstack[c]{\strut #1\strut}}

% ---------- Lines ----------
\newcommand{\TitleLine}{\textcolor{StageBlue}{\rule{0.98\linewidth}{0.95pt}}}
\newcommand{\SubTitleLine}{\textcolor{RuleGray}{\rule{0.98\linewidth}{0.65pt}}}

% ---------- Continued on next page ----------
\newcommand{\ContinuedOnNextPageInline}{%
  \par
  \vfill
  \noindent\textcolor{RuleGray}{\rule{\linewidth}{1pt}}\vspace{-2pt}
  \noindent\hfill\textit{Continued on next page}\par\vspace{-10pt}
  \noindent\textcolor{RuleGray}{\rule{\linewidth}{1pt}}
  \clearpage
}

% ---------- PDF hyperlinks ----------
\usepackage[hidelinks]{hyperref}
\hypersetup{
  colorlinks=false,
  pdfborder={0 0 0},
  linktoc=all,
  pdftitle={\DocTitle},
  pdfauthor={\ProgramName}
}

% =========================================================
\begin{document}

\subsection*{Phase 03 Card Header}
\phantomsection
\addcontentsline{toc}{subsection}{Phase 03 Card Header}

\begin{itemize}
\item Phase \textbf{ID}: Phase 03
\item Phase Name: Converting Strength to Power and Endurance for SOF Selection
\item Duration weeks: 12
\item Version: v1.0
\item Governing Status: Deployable
\end{itemize}

\subsection*{1) Phase Mission}
\phantomsection

Convert Phase 02 strength reserves into repeatable endurance outputs and higher specificity \textbf{RUN} performance while maintaining \textbf{CAL} standardization and strict validity discipline for gate-grade evidence. Phase 03 must increase selection-relevant output capacity without permitting unsafe stacking, informal test substitution, or durability breakdown that would make evidence non-repeatable or unsafe to reproduce.

Mission boundaries (explicit):

\begin{itemize}
\item Phase 03 is not permitted to drift into unscheduled, maximal testing behavior during build weeks.
\item Phase 03 does not authorize “count it anyway” evidence; \textbf{INVALID} tests are treated as not yet tested for gate decisions.
\item Phase 03 does not authorize replacing \textbf{RUN} gate constructs with non-run conditioning outputs.
\item Phase 03 increases \textbf{RUN} specificity and endurance verification only under protected-window conditions and under durability constraints.
\end{itemize}

\subsection*{2) Phase Purpose}
\phantomsection

Phase 03 exists to translate the improved strength ceiling from Phase 02 into usable, repeatable work output under fatigue constraints, with specific emphasis on:

\begin{itemize}
\item \textbf{RUN} verification expansion: transitioning toward longer verification distances (\textbf{C10} / \textbf{MET-2B-016}) under validity-protected conditions.
\item \textbf{CAL} maintenance with improved repeatability: sustain strict execution rules while reducing technique drift and preventing counting ambiguity.
\item Durability stabilization as a constraint, not a “side goal”: gait stability and foot integrity determine whether higher-cost tests are eligible.
\item Evidence integrity consolidation: Phase 03 ensures that test events are comparable across windows by enforcing stable environment identity, stable tool identity, stable footwear interface, and complete logging.
\end{itemize}

Not the purpose (explicit exclusions of intent drift):

\begin{itemize}
\item Not to invent new tests (no interval benchmarks, no mixed-modal simulations) unless they already exist as Stage 2 protocolized instruments.
\item Not to compress missed verification into stacked weeks.
\item Not to pursue “peak day” performance claims; replication and admissibility govern evidence use.
\end{itemize}

\subsection*{3) Entry Criteria}
\phantomsection

Entry criteria are objective and auditable. They are satisfied only when the transition from Phase 02 to Phase 03 can be justified using admissible evidence and stable risk posture.

A) Prior gate outcome requirement

\begin{itemize}
\item Phase 02 must conclude with an \textbf{ADVANCE} decision recorded under Stage 1.F gate doctrine.
\end{itemize}

B) Decision window compliance posture (Stage 1.F default eligibility)

\begin{itemize}
\item The two most recent complete training weeks preceding Phase 03 Week 1 must demonstrate:
  \begin{itemize}
  \item at least 5 of 6 admissible sessions per week on average,
  \item outside medical issues, no more than 1 \textbf{MODIFIED} session per week admissible,
  \item zero \textbf{INVALID} sessions in the decision window (binary disqualifier posture for \textbf{ADVANCE}).
  \end{itemize}
\end{itemize}

C) Durability prerequisites (disqualifier controls)

\begin{itemize}
\item \textbf{MET-2B-036} (Gait-Altering Pain Flag) is No across the decision window.
\item \textbf{MET-2B-037} (Foot Hotspot or Blister Flag) does not show an active breakdown pattern that would predictably compromise \textbf{RUN} mechanics during Phase 03 verification.
\end{itemize}

D) Test feasibility prerequisites (instrument readiness posture)

\begin{itemize}
\item \textbf{CAL} artifact workflow is executable:
  \begin{itemize}
  \item camera placement and angle supports auditable rep criteria,
  \item timer visibility/audibility supports auditable start/stop,
  \item Evidence Ref workflow is stable (every test event can be binder-referenced).
  \end{itemize}
\item \textbf{RUN} environment/tool identity is executable:
  \begin{itemize}
  \item either a stable, measurable route identity exists, or
  \item a stable treadmill identity exists and the required settings can be recorded per Stage 2 validity posture.
  \end{itemize}
\end{itemize}

E) Default posture if any entry criterion is not satisfied

\begin{itemize}
\item Default posture is \textbf{HOLD} at phase transition, with corrective action and reslot per Stage 3.E. Phase 03 gate escalation posture must not begin from an unstable, inadmissible, or high-risk entry state.
\end{itemize}

\subsection*{4) Operating Constraints}
\phantomsection

Operating constraints are non-negotiable and apply to every week in Phase 03.

A) Cadence and session architecture

\begin{itemize}
\item Weekly cadence is fixed: 6 sessions per week, no two-a-days, no make-up doubling.
\item Every session includes the required elements in the required order per Stage 1.F:
  \begin{itemize}
  \item Safety self-check first
  \item Warm-up
  \item Mobility
  \item Main work
  \item Cardio
  \item Cooldown
  \item Recovery notes and post-session notes
  \end{itemize}
\end{itemize}

B) Cardio minimum and continuity standard (body text enforcement)

\begin{itemize}
\item Cardio occurs in every training session.
\item Cardio minimum standard is minimum 30 minutes continuous per training session.
\item If Cardio is under the minimum or not continuous, the session is \textbf{INVALID} and provides no gate credit.
\end{itemize}

C) Daily baseline activity

\begin{itemize}
\item 10,000 steps per day is a baseline activity requirement and must be logged daily.
\item Steps do not count as a session and do not replace session structure requirements.
\end{itemize}

D) Evidence and admissibility posture

\begin{itemize}
\item Gate decisions rely on admissible evidence only:
  \begin{itemize}
  \item sessions must be \textbf{VALID} to count toward gate eligibility (\textbf{MODIFIED} admissible only within Stage 1.F cap posture),
  \item tests must be tagged \textbf{VALID} or \textbf{INVALID} only,
  \item \textbf{INVALID} tests provide no gate credit and default to reslot posture.
  \end{itemize}
\end{itemize}

E) High-level scope boundaries (Phase 03 execution posture)

\begin{itemize}
\item Gate-grade tests are executed only in protected deload/test windows defined in Section 9.
\item Safety overrides schedule and performance.
\item Any attempt to bypass stop posture or force testing under compromised conditions contaminates evidence and defaults gate posture conservative.
\end{itemize}

\subsection*{5) Primary Adaptations and Physiological Focus}
\phantomsection

Phase 03 adaptation targets are expressed as governed outcomes (what must change) without programming detail.

\begin{itemize}
\item \textbf{RUN} endurance verification capacity: improved ability to produce repeatable performance outcomes at longer distance (3-mile) under recorded conditions and stable mechanics.
\item Power-endurance conversion at a systems level: improved ability to sustain output without acute collapse, while maintaining adherence and evidence quality.
\item \textbf{CAL} endurance repeatability: improved stability of strict standards under fatigue and improved consistency of outcomes across protected windows.
\item Durability as a gating control: reduced incidence of gait-altering pain flags and foot breakdown, enabling lawful verification without invalidity.
\item Recovery stability: stable sleep and fatigue posture sufficient to avoid systematic fatigue contamination of test windows.
\item Compliance stability: reduced frequency of \textbf{MODIFIED}/\textbf{INVALID} sessions, stable evidence capture, and stable environment/tool identity posture.
\end{itemize}

\subsection*{6) Primary Modalities}
\phantomsection

Primary domains emphasized in Phase 03 (Stage 2 tags):

\begin{itemize}
\item \textbf{RUN}: increased verification emphasis; end-phase gate uses 3-mile construct.
\item \textbf{CAL}: continued standardized battery as a stable evidence stream; gate-grade evidence requires compliant artifacts.
\item \textbf{DURABILITY}: ongoing screen and constraint posture; disqualifiers constrain gate eligibility and test execution.
\end{itemize}

Secondary support domains:

\begin{itemize}
\item \textbf{STR}: maintenance posture supporting conversion without dominating fatigue budget.
\item \textbf{RECOVERY}: readiness markers used to determine whether a protected window is test-capable or must convert to deload-only/check-in-only.
\item \textbf{BODYCOMP}: trend continuity supports feasibility and durability; it does not substitute for performance verification.
\end{itemize}

Explicit non-equivalency doctrine (Phase 03):

\begin{itemize}
\item \textbf{AER} constructs are not used as evidence for \textbf{RUN} performance or \textbf{RUN} tolerance.
\item Non-standard “interval benchmarks” are not scheduled or used as gate evidence unless they exist as Stage 2 protocolized tests.
\end{itemize}

\subsection*{7) Weekly Structure (Sessions 1–6)}
\phantomsection

The weekly structure is intent-level only. It defines the role of each session in maintaining Phase 03 intent while controlling fatigue and preserving validity.

\clearpage
\begin{landscape}

\begingroup
\scriptsize
\setlength{\tabcolsep}{2pt}
\renewcommand{\arraystretch}{1.12}

\begin{longtblr}{
  width=\linewidth,
  rowhead=1,
  colspec = {
    Q[l,wd=0.70in]
    X[l,1.65]
    X[l,1.15]
    X[l,2.55]
    X[l,2.85]
  },
  hlines,
  vlines,
  cells = {hsep=6pt, vsep=6pt, halign=l, valign=t},
  cell{1-Z}{1-5} = {fg=black},
  row{odd} = {bg=white},
  row{even} = {bg=LightGray},
  row{1} = {bg=StageBlue!15, fg=black, font=\bfseries, halign=l, valign=t},
}
\Hdr{Session \#} &
\Hdr{Primary Focus} &
\Hdr{Main Work Domains} &
\Hdr{Cardio Modality/Intent} &
\Hdr{Notes/Risk Controls} \\

1 &
\textbf{RUN} mechanics, economy, and stable throughput &
\textbf{RUN}; \textbf{DURABILITY}; \textbf{RECOVERY} &
Modality: run or treadmill run (as tolerated); Minutes: 30–60; Intensity Anchor: RPE 3–5 with talk-test full-to-short sentences; Continuity: continuous planned &
This session is not a time trial. If \textbf{MET-2B-036} becomes Yes or foot hotspots escalate, substitute to lower-impact Cardio and document; do not push through mechanics change. No novelty near protected windows. \\

2 &
\textbf{CAL} standardization and trunk endurance &
\textbf{CAL}; \textbf{DURABILITY} &
Modality: low-impact steady (walk/bike/elliptical); Minutes: 30–45; Intensity Anchor: RPE 2–4 full sentences; Continuity: continuous planned &
Maintain strict \textbf{CAL} rules in all practice exposures; reserve gate-grade efforts for protected windows with artifacts. Protect elbows/shoulders/wrists; terminate on pain-driven compensation. \\

3 &
Aerobic support and recovery throughput &
\textbf{RECOVERY}; \textbf{DURABILITY} &
Modality: lowest-risk continuous Cardio; Minutes: 30–50; Intensity Anchor: RPE 2–4 full sentences; Continuity: continuous planned &
Acts as fatigue-buffer day to protect later high-cost sessions and preserve compliance. If recovery is degraded, this becomes Minimum Viable Session posture while preserving Cardio continuity and logging. \\

4 &
\textbf{STR} maintenance and movement quality &
\textbf{STR}; \textbf{DURABILITY} &
Modality: low-impact steady; Minutes: 30–45; Intensity Anchor: RPE 2–4 full sentences; Continuity: continuous planned &
\textbf{STR} does not own gate tests in this phase. Avoid fatigue spikes that would contaminate \textbf{RUN}/\textbf{CAL} protected windows. Any mechanics change triggers regression and conservative downshift. \\

5 &
\textbf{RUN} endurance support without time-trial load &
\textbf{RUN}; \textbf{RECOVERY}; \textbf{DURABILITY} &
Modality: run or treadmill run, or substitute when constrained; Minutes: 30–60; Intensity Anchor: RPE 3–5; Talk Test: full-to-short sentences; Continuity: continuous &
Avoid mid-week “prove-it” behavior. Preserve stable route/treadmill identity. If weather hazards exist, substitute indoors and document; do not force unsafe outdoor exposure for compliance optics. \\

6 &
Consolidation and pre-window stabilization &
\textbf{CAL} (light) or \textbf{RUN} (light); \textbf{RECOVERY} &
Modality: low-impact steady; Minutes: 30–45; Intensity Anchor: RPE 2–4 full sentences; Continuity: continuous planned &
Stabilization day, especially before Week 4/8/12 windows. Confirm artifacts, timing method, route/treadmill identity, and logging readiness. No novelty, no intensity spikes, no catch-up behavior. \\

\end{longtblr}

\endgroup

\end{landscape}
\clearpage

\subsection*{8) Progression Rules (Volume and Intensity Caps; Regression Triggers)}
\phantomsection

Progression is governed by stability and admissibility, not by motivation or isolated good days.

A) One-lever change rule (global)

In a single week, change only one of the following levers:

\begin{itemize}
\item total weekly Cardio minutes beyond the minimum,
\item longest Cardio-primary session duration,
\item \textbf{RUN}-specific exposure complexity (without moving into unscheduled test efforts),
\item \textbf{CAL} exposure density (without turning practice into maximal attempts).
\end{itemize}

B) \textbf{RUN} progression governance

\begin{itemize}
\item \textbf{RUN} is progressed only when:
  \begin{itemize}
  \item \textbf{MET-2B-036} remains No,
  \item foot interface is stable (no active breakdown pattern),
  \item recovery markers do not indicate fatigue contamination risk approaching protected windows.
  \end{itemize}
\item If mechanics change occurs:
  \begin{itemize}
  \item terminate provoking exposure immediately,
  \item substitute to lower-risk modality,
  \item treat any planned gate-grade test in the next window as at-risk for invalidity and apply reslot posture if stability is not restored.
  \end{itemize}
\end{itemize}

C) \textbf{CAL} progression governance

\begin{itemize}
\item \textbf{CAL} practice is technique-driven and standards-driven:
  \begin{itemize}
  \item do not chase maximal attempts outside protected windows,
  \item treat the artifact workflow as part of validity infrastructure (practice the setup without necessarily executing maximal efforts).
  \end{itemize}
\item Gate-grade \textbf{CAL} evidence is produced only in protected windows and only with compliant artifacts; otherwise it is \textbf{INVALID} for gate use.
\end{itemize}

D) Regression triggers (same-day and trend-based)

Same-day triggers (terminate and regress):

\begin{itemize}
\item pain that alters mechanics,
\item dizziness/lightheadedness,
\item unusual/disproportionate breathlessness relative to planned effort,
\item instability/swelling,
\item stop-rule symptoms per Stage 0.5 / Stage 1.F.
\end{itemize}

Trend triggers (convert protected window to deload-only or check-in-only and reslot maximal tests):

\begin{itemize}
\item sleep debt trend and rising session RPE at stable nominal workload,
\item escalating foot hotspots or early blister signs,
\item repeated \textbf{MODIFIED} / \textbf{INVALID} sessions in the lead-in to a protected window,
\item environment or tool drift not stabilized (route changes, treadmill changes, timing method changes).
\end{itemize}

\subsection*{9) Deload and Test Placement (Mid-Phase and End-of-Phase)}
\phantomsection

Phase 03 uses protected windows to produce gate-grade evidence and to protect test validity from fatigue contamination.

Protected deload/test windows (week-in-phase; scheduling doctrine):

\begin{itemize}
\item Week 4: Deload+Test mid-phase touchpoint
\item Week 8: Deload+Test mid-phase touchpoint
\item Week 12: Deload+Test end-phase exit gate
\end{itemize}

Protected window rules:

\begin{itemize}
\item Gate-grade tests occur only in protected windows.
\item Major test events must respect Stage 3 anti-stacking and spacing intent.
\item If validity cannot be preserved due to environment, tool drift, illness, or durability flags, the protected window converts to deload-only/check-in-only and the affected major tests are reslotted to the next protected window.
\end{itemize}

\subsection*{10) Test Battery (IDs and Labels Only)}
\phantomsection

Test battery entries are identifiers only and do not reproduce protocol procedure text.

\subsubsection*{Mid-Phase Tests (Week 4 and Week 8)}
\phantomsection

\begin{itemize}
\item \textbf{CAL} battery:
  \begin{itemize}
  \item \textbf{C4} / \textbf{MET-2B-010} Push-Ups — 2-Minute Timed, Strict Standard
  \item \textbf{C5} / \textbf{MET-2B-011} Sit-Ups — 2-Minute Timed, Standard
  \item \textbf{C6} / \textbf{MET-2B-012} Pull-Ups — Strict Dead-Hang
  \item \textbf{C7} / \textbf{MET-2B-013} Plank — Forearm
  \end{itemize}
\item \textbf{RUN} check-in:
  \begin{itemize}
  \item Week 4: \textbf{C8} / \textbf{MET-2B-014} 1.5-Mile Run Time Trial (check-in construct)
  \item Week 8: \textbf{C9} / \textbf{MET-2B-015} 2-Mile Run Time Trial (check-in construct)
  \end{itemize}
\item Durability readiness screen (low-cost; may be recorded within the same window without counting as a major event):
  \begin{itemize}
  \item \textbf{C16} / \textbf{MET-2B-035} Pain Severity; \textbf{MET-2B-036} Gait-Altering Pain Flag; \textbf{MET-2B-037} Foot Hotspot or Blister Flag
  \end{itemize}
\end{itemize}

\subsubsection*{End-Phase Tests (Week 12)}
\phantomsection

\begin{itemize}
\item \textbf{CAL} battery:
  \begin{itemize}
  \item \textbf{C4} / \textbf{MET-2B-010}
  \item \textbf{C5} / \textbf{MET-2B-011}
  \item \textbf{C6} / \textbf{MET-2B-012}
  \item \textbf{C7} / \textbf{MET-2B-013}
  \end{itemize}
\item \textbf{RUN} exit gate construct:
  \begin{itemize}
  \item \textbf{C10} / \textbf{MET-2B-016} 3-Mile Run Time Trial
  \end{itemize}
\item Durability readiness screen:
  \begin{itemize}
  \item \textbf{C16} / \textbf{MET-2B-035}; \textbf{MET-2B-036}; \textbf{MET-2B-037}
  \end{itemize}
\end{itemize}

\begin{landscape}

\subsubsection*{Test Battery Table}
\phantomsection

\begingroup
\scriptsize
\setlength{\tabcolsep}{2pt}
\renewcommand{\arraystretch}{1.12}

\begin{longtblr}{
  width=\linewidth,
  rowhead=1,
  colspec = {
    Q[l,wd=0.95in]
    X[l,2.35]
    Q[l,wd=0.95in]
    X[l,2.65]
    X[l,1.35]
  },
  hlines,
  vlines,
  cells = {hsep=6pt, vsep=6pt, halign=l, valign=t},
  cell{1-Z}{1-5} = {fg=black},
  row{odd} = {bg=white},
  row{even} = {bg=LightGray},
  row{1} = {bg=StageBlue!15, fg=black, font=\bfseries, halign=l, valign=t},
}
\Hdr{Test Timing} &
\Hdr{Test \textbf{ID}/Name} &
\Hdr{Domain} &
\Hdr{Validity Conditions} &
\Hdr{Pass Threshold} \\

Week 4 &
\textbf{C16} / \textbf{MET-2B-035}; \textbf{MET-2B-036}; \textbf{MET-2B-037} Durability and Readiness Screen &
\textbf{DURABILITY} &
Required fields complete; same-day capture; no diagnostic language; gait-altering pain flag constrains test eligibility posture and may convert window to check-in &
Pass when recorded and disqualifiers are not active \\

Week 4 &
\textbf{C4} / \textbf{MET-2B-010}; \textbf{C5} / \textbf{MET-2B-011}; \textbf{C6} / \textbf{MET-2B-012}; \textbf{C7} / \textbf{MET-2B-013} \textbf{CAL} battery &
\textbf{CAL} &
\textbf{CAL} artifact evidence compliant; warm-up recorded; safety screen passed; complete required fields; Evidence Ref recorded &
Tier per Stage 2 for check-in (no gate decision on check-in alone) \\

Week 4 &
\textbf{C8} / \textbf{MET-2B-014} 1.5-Mile Run Time Trial &
\textbf{RUN} &
Route or treadmill identity logged; warm-up recorded; safety screen passed; no stop-rule violations; Evidence Ref recorded &
Tier per Stage 2 for check-in \\

Week 8 &
\textbf{C16} / \textbf{MET-2B-035}; \textbf{MET-2B-036}; \textbf{MET-2B-037} Durability and Readiness Screen &
\textbf{DURABILITY} &
Required fields complete; gait-altering pain flag constrains gate posture &
Pass when recorded and disqualifiers are not active \\

Week 8 &
\textbf{C4} / \textbf{MET-2B-010}; \textbf{C5} / \textbf{MET-2B-011}; \textbf{C6} / \textbf{MET-2B-012}; \textbf{C7} / \textbf{MET-2B-013} \textbf{CAL} battery &
\textbf{CAL} &
\textbf{CAL} artifact evidence compliant; complete required fields; Evidence Ref recorded &
Tier per Stage 2 for mid-phase gate focus \\

Week 8 &
\textbf{C9} / \textbf{MET-2B-015} 2-Mile Run Time Trial &
\textbf{RUN} &
Route or treadmill identity logged; required settings logged when treadmill is used; no unsafe environment exposure; Evidence Ref recorded &
Tier per Stage 2 for mid-phase gate focus \\

Week 12 &
\textbf{C16} / \textbf{MET-2B-035}; \textbf{MET-2B-036}; \textbf{MET-2B-037} Durability and Readiness Screen &
\textbf{DURABILITY} &
Must be \textbf{VALID}; gait-altering pain flag must be No for advancement posture &
Pass if gait-altering pain flag is No and foot breakdown not active \\

Week 12 &
\textbf{C4} / \textbf{MET-2B-010}; \textbf{C5} / \textbf{MET-2B-011}; \textbf{C6} / \textbf{MET-2B-012}; \textbf{C7} / \textbf{MET-2B-013} \textbf{CAL} battery &
\textbf{CAL} &
Must be \textbf{VALID}; compliant artifacts; Evidence Ref recorded; replication posture applies for tier classification used in gates &
\textbf{INTERMEDIATE} tier per Stage 2 (with replication) \\

Week 12 &
\textbf{C10} / \textbf{MET-2B-016} 3-Mile Run Time Trial &
\textbf{RUN} &
Must be \textbf{VALID}; route or treadmill identity logged; measurement integrity satisfied; Evidence Ref recorded; replication posture applies for tier classification used in gates &
\textbf{INTERMEDIATE} tier per Stage 2 (with replication) \\

\end{longtblr}

\endgroup

\end{landscape}
\clearpage

\subsection*{11) Exit Standards (Objective Pass and Fail)}
\phantomsection

Phase 03 exit gate is evaluated in Week 12 protected window only and uses admissible evidence only.

A) Eligibility to decide (Stage 1.F decision window posture)

\begin{itemize}
\item Decision window: the two most recent complete training weeks preceding Week 12.
\item Minimum admissible session posture: at least 5 of 6 admissible sessions per week on average.
\item \textbf{MODIFIED} admissibility cap: outside medical issues, no more than 1 \textbf{MODIFIED} session per week admissible.
\item Disqualifier: any \textbf{INVALID} session within the decision window disqualifies \textbf{ADVANCE} and forces \textbf{HOLD}.
\end{itemize}

B) Required gate-grade tests (Week 12; \textbf{VALID} only)

Pass conditions:

\begin{itemize}
\item \textbf{RUN}:
  \begin{itemize}
  \item \textbf{C10} / \textbf{MET-2B-016} is \textbf{VALID} and meets \textbf{INTERMEDIATE} tier posture per Stage 2.
  \end{itemize}
\item \textbf{CAL}:
  \begin{itemize}
  \item \textbf{C4} / \textbf{MET-2B-010}, \textbf{C5} / \textbf{MET-2B-011}, \textbf{C6} / \textbf{MET-2B-012}, \textbf{C7} / \textbf{MET-2B-013} are \textbf{VALID} and meet \textbf{INTERMEDIATE} tier posture per Stage 2, including compliant artifacts.
  \end{itemize}
\item Durability constraints:
  \begin{itemize}
  \item \textbf{MET-2B-036} must be No in the decision window and at the time of testing.
  \item \textbf{MET-2B-037} must not indicate an active breakdown state that would make outputs non-repeatable or unsafe to reproduce.
  \end{itemize}
\end{itemize}

Fail conditions:

\begin{itemize}
\item Any required gate test is \textbf{INVALID}.
\item Any required gate test lacks Evidence Ref, required logging fields, or required artifacts (\textbf{CAL}).
\item Any stop-rule event occurs during a gate test or is active and unresolved.
\item Any gait-altering pain flag is Yes in the decision window.
\end{itemize}

C) Replication posture (tier use)

\begin{itemize}
\item Tier classification used for advancement requires the Stage 2 replication rule (2 of last 3 \textbf{VALID} test events within the decision window).
\item If replication is not met for a gate-critical metric, default posture is \textbf{HOLD} for that metric until replication exists.
\end{itemize}

\subsection*{12) Risk Controls and Stop Rules (Phase-Specific Overlays)}
\phantomsection

Stage 0.5 and Stage 1.F remain controlling. Phase 03 overlays apply them with additional specificity for higher-cost endurance verification.

\textbf{RUN}-specific overlays:

\begin{itemize}
\item No gate-grade \textbf{RUN} attempts outside protected windows.
\item Route/tool stability is treated as a validity prerequisite:
  \begin{itemize}
  \item route identity or treadmill identity must remain stable in the lead-in to Week 12,
  \item novelty near the gate window is treated as a comparability break and may require reslot posture.
  \end{itemize}
\item Environment hazard handling is mandatory:
  \begin{itemize}
  \item unsafe footing, extreme cold/wind, poor visibility, unsafe heat, poor air quality triggers indoor substitution only when Stage 2 validity and logging preserve comparability; otherwise reslot.
  \end{itemize}
\end{itemize}

\textbf{CAL}-specific overlays:

\begin{itemize}
\item \textbf{CAL} gate-grade evidence requires compliant artifacts; missing/non-compliant artifacts default tests to \textbf{INVALID} for gate use.
\item Avoid fatigue stacking:
  \begin{itemize}
  \item do not schedule other major events in the same protected window in a way that compromises \textbf{CAL} validity (e.g., performing \textbf{CAL} battery after an uncontrolled high-fatigue day).
  \end{itemize}
\end{itemize}

Durability overlays (hard constraints):

\begin{itemize}
\item \textbf{MET-2B-036} Yes is a disqualifier signal that blocks safe escalation and invalidates claims of readiness for advancement posture until resolved.
\item Foot integrity is treated as a limiting system; recurring hotspot escalation triggers downshift and reslot of \textbf{RUN} gate attempts.
\end{itemize}

Stop posture integration:

\begin{itemize}
\item \textbf{STOP} \textbf{NOW} triggers terminate the attempt and convert the test to \textbf{INVALID} with the appropriate \textbf{TEST-RC}. The corrective action is reslot, not “finish anyway.”
\end{itemize}

\subsection*{13) Required Capabilities and Dependencies}
\phantomsection

Dependencies must be validated prior to gate attempts; if not validated, default posture is \textbf{HOLD} and reslot per Stage 3.

\clearpage
\begin{landscape}

\begingroup
\scriptsize
\setlength{\tabcolsep}{2pt}
\renewcommand{\arraystretch}{1.12}

\begin{longtblr}{
  width=\linewidth,
  rowhead=1,
  colspec = {
    Q[l,wd=0.95in]
    X[l,2.55]
    Q[l,wd=1.15in]
    Q[l,wd=1.05in]
    X[l,3.10]
  },
  hlines,
  vlines,
  cells = {hsep=6pt, vsep=6pt, halign=l, valign=t},
  cell{1-Z}{1-5} = {fg=black},
  row{odd} = {bg=white},
  row{even} = {bg=LightGray},
  row{1} = {bg=StageBlue!15, fg=black, font=\bfseries, halign=l, valign=t},
}
\Hdr{Dependency \textbf{ID}} &
\Hdr{Required Capability} &
\Hdr{Due-By} &
\Hdr{Owner} &
\Hdr{Fail Action} \\

\textbf{DEP-1C-014} &
7.08 Conditioning, Endurance, and Aerobic Training Tools — minimum viable safe \textbf{RUN} environment (route or treadmill) + substitution capability under environment volatility &
Before Week 4 &
Trainee &
If safe route or indoor substitute is not available, \textbf{RUN} tests are not executed as gate-grade; reslot to next protected window; preserve Cardio via lowest-risk substitutes; default \textbf{HOLD} if end gate cannot be executed validly. \\

\textbf{DEP-1C-019} &
7.14 Environmental Apparel — minimum capability to execute safely outdoors or shift indoors; environment go/no-go rules are written and used &
Before Week 4 &
Coach Lead &
If environment plan is missing or not followed, any compromised tests are \textbf{INVALID}; reslot; prohibit outdoor testing under hazard conditions; default \textbf{HOLD} until execution is controlled. \\

\textbf{DEP-1C-021} &
7.16 Operational Self-Management Systems — same-day logging capability and compliance proof &
Before Week 4 &
Trainee &
Missing same-day logs make evidence inadmissible; \textbf{HOLD} for gate decisions; repeat compliance proof; do not attempt gate-grade tests until log discipline is stable. \\

\textbf{DEP-1C-006} &
7.18 Testing, Timing, Measurement, and Course-Setup Tools — timer method stability; tool identity loggable; route/treadmill identity loggable &
Before Week 4 &
Equipment Lead &
If timing or tool identity cannot be recorded, \textbf{RUN} tests are \textbf{INVALID} for gate use; reslot; no advancement on missing measurement integrity. \\

\textbf{DEP-1C-007} &
7.11 Footwear Systems and Foot Care — footwear interface stable and verified; no novelty near protected windows &
Before Week 4 and maintained through Week 12 &
Equipment Lead &
If hotspots/blisters trend upward or gait is altered, downshift exposure and reslot \textbf{RUN} gate attempts; \textbf{HOLD} posture for advancement until foot integrity stabilizes. \\

\textbf{DEP-1C-008} &
7.11 Footwear Systems and Foot Care — footcare protocol supplies and immediate hotspot response posture &
Before Week 4 &
Equipment Lead &
If hotspot protocol cannot be executed, cap step-based exposure and preserve Cardio via lower-friction modalities; reslot gait-dependent tests; default \textbf{HOLD} until protocol is stable. \\

\textbf{DEP-1C-021} &
7.16 Operational Self-Management Systems — schedule stability for 6 sessions/week and decision window integrity &
By Week 2 &
Trainee &
If cadence collapses and decision window becomes non-admissible, default \textbf{HOLD}; re-establish cadence proof before attempting end gate. \\

\textbf{DEP-1C-006} &
7.18 Testing, Timing, Measurement, and Course-Setup Tools — test environment measurability and environment logging capability &
Before Week 8 &
Coach Lead &
Tests under unverified conditions are \textbf{INVALID} for gate use; reslot; prohibit “estimated” distances for gate-grade \textbf{RUN} tests. \\

\textbf{DEP-1C-016} &
7.10 Logistics, Maintenance, Storage, and Sustainment Equipment — drying and sustainment posture sufficient to prevent repeatable foot breakdown &
By Week 2 &
Trainee &
If drying/sustainment fails and foot breakdown recurs, downshift and reslot; treat as durability constraint; \textbf{HOLD} advancement until stabilized. \\

\textbf{DEP-1C-013} &
7.13 Nutrition Preparation and Hydration Systems — hydration access and container availability; deviations are loggable &
Before Week 4 &
Trainee &
If hydration instability produces safety ambiguity, downshift intensity, avoid heat exposure, and reslot gate tests; default \textbf{HOLD} if safety and validity cannot be preserved. \\

\end{longtblr}

\endgroup

\end{landscape}
\clearpage

\subsection*{14) Validity Requirements (Evidence Admissibility)}
\phantomsection

Validity requirements are enforced controls; they determine whether evidence can be used for gate decisions.

Session admissibility requirements:

\begin{itemize}
\item Sessions must include all required elements and minimum dataset per Stage 1.F.
\item Cardio is minimum 30 minutes continuous per session; failure makes the session \textbf{INVALID}.
\item \textbf{SESSION-RC} codes are used when sessions are \textbf{MODIFIED} or \textbf{INVALID}; undocumented substitutions are treated as \textbf{INVALID}.
\end{itemize}

Test admissibility requirements:

\begin{itemize}
\item Tests must be tagged \textbf{VALID} or \textbf{INVALID} only.
\item Evidence Ref is required only when the governing Stage 2 protocol or logging rule requires it; missing Evidence Ref defaults the test to \textbf{INVALID} for gate use.
\item \textbf{CAL} tests used for gate decisions require compliant artifact evidence; missing or non-compliant artifact defaults to \textbf{INVALID} for gate use.
\item \textbf{RUN} tests used for gate decisions require:
  \begin{itemize}
  \item stable route identity or treadmill identity,
  \item required settings logging where applicable,
  \item warm-up completion record,
  \item stop posture honored,
  \item environment hazards logged and handled via substitution/reslot discipline.
  \end{itemize}
\end{itemize}

Replication requirement for gate use:

\begin{itemize}
\item Tier classification used for advancement requires 2 of last 3 \textbf{VALID} test events within the decision window for that metric.
\item If replication is not met, the result may be recorded, but the metric posture defaults \textbf{HOLD} for progression decisions.
\end{itemize}

\subsection*{15) Change Control Notes (Freeze Windows and Patch Discipline)}
\phantomsection

Freeze and change-control posture is applied to protect comparability.

Allowed without patch (documentation required; must not alter stage doctrine):

\begin{itemize}
\item Swap the sequence of tests inside the same protected window week to preserve spacing and safety, without adding tests and without weakening validity requirements.
\item Reslot a missed/invalid test to the next protected window within the phase.
\item Convert a protected window into deload-only or check-in-only when readiness, durability, or environment constraints make gate-grade testing inadmissible.
\item Substitute \textbf{RUN} route to treadmill only when Stage 2 validity/logging preserves comparability; otherwise reslot.
\end{itemize}

Requires patch (change-control; revalidation required):

\begin{itemize}
\item Changing the gate test set (changing \textbf{RUN} distance construct or altering \textbf{CAL} battery composition beyond Stage 2).
\item Changing protected window placement or moving gate decisions into build weeks.
\item Extending phase duration or adding new protected windows.
\item Any change that alters admissibility rules, required logging fields, or artifact requirements.
\end{itemize}

\subsection*{16) Compliance Checklist (Pass and Fail)}
\phantomsection

Pass conditions (minimum):

\begin{itemize}
\item Phase identity fields complete: Phase 03, 12 weeks, version v1.0, status Deployable.
\item Weekly structure table includes Sessions 1–6 and includes Cardio Modality/Intent entries with modality, minutes, RPE + talk-test cue, and continuity statement.
\item Operating constraints explicitly state:
  \begin{itemize}
  \item 6 sessions/week,
  \item required session elements in correct order,
  \item Cardio minimum 30 minutes continuous per session in body text,
  \item 10,000 steps/day baseline in body text,
  \item validity/admissibility posture (\textbf{VALID}-only gate evidence; \textbf{INVALID} provides no gate credit).
  \end{itemize}
\item Protected windows are explicit (Week 4, Week 8, Week 12) and gate-grade tests are stated as occurring only within those windows.
\item Test Battery references Stage 2 instruments only and includes no placeholders:
  \begin{itemize}
  \item \textbf{CAL}: \textbf{C4}–\textbf{C7} with \textbf{MET-2B-010}–013
  \item \textbf{RUN}: \textbf{C8}–\textbf{C10} with \textbf{MET-2B-014}–016
  \item \textbf{DURABILITY}: \textbf{C16} with \textbf{MET-2B-035}–037
  \item Compliance: \textbf{MET-2B-044}–046
  \end{itemize}
\item Dependencies table references \textbf{DEP-1C}-\#\#\# and uses Stage 7 category IDs 7.01–7.20 only, does not require 7.02, and includes owners and explicit fail actions defaulting to \textbf{HOLD}/reslot rather than unsafe forcing.
\item Exit standards are objective, tier-referenced (per Stage 2), and enforce replication posture and disqualifier posture.
\end{itemize}

Fail conditions (non-deployable triggers):

\begin{itemize}
\item Any gate test listed without Stage 2 identifiers.
\item Any implication that gate-grade evidence is produced in build weeks.
\item Any use of non-equivalency violations (e.g., treating non-run conditioning as \textbf{RUN} gate evidence).
\item Any missing dependency owner, missing fail action, or missing protected window placement posture that would allow uncontrolled drift.
\end{itemize}

\end{document}